%%%%%%%%%%%%%%%%%%%%%%%%%%%%%%%%%%%%%%%%%
% Otto-Friedrich-University of Bamberg
% Chair of Explainable Machine Learning (xAI)
% Deep Learning Assignments
% LaTeX Template
% Version 1.0
%
% Original template has been downloaded from:
% https://de.overleaf.com/latex/templates/maths-coursework-template/kbyhcwmjdtpf
%
% Original author:
% Qiao Han
%
% License:
% Creative Commons CC BY 4.0 (https://creativecommons.org/licenses/by/4.0/)
%
% Author:
% Sebastian Dörrich
%%%%%%%%%%%%%%%%%%%%%%%%%%%%%%%%%%%%%%%%%
%
%----------------------------------------------------------------------------------------
% USED PACKAGES AND OTHER DOCUMENT CONFIGURATIONS
%----------------------------------------------------------------------------------------
\usepackage{ifxetex}
\usepackage{textpos}
\usepackage[round]{natbib} % \citet{key} -> Author et al. (Year); papers live in references.bib
\usepackage{kpfonts}
\usepackage[a4paper,hmargin=2.8cm,vmargin=2.0cm,includeheadfoot]{geometry}
\usepackage{ifxetex}
\usepackage{stackengine}
\usepackage{tabularx,longtable,multirow,caption, subcaption} % hangcaption
\usepackage{fncylab} % formatting of labels
\usepackage{fancyhdr}
%\usepackage{color}
\usepackage[tight,ugly]{units}
\usepackage{url}
\usepackage{float}
\usepackage[english]{babel}
\usepackage{amsmath}
\usepackage{graphicx}
% todonotes was loaded but never used, and its \todo clashed with the todo
% environment defined below. Removed; re-add it if you ever want margin notes.
%\usepackage[colorinlistoftodos]{todonotes}
\usepackage{dsfont}
\usepackage{epstopdf} % automatically replace .eps with .pdf in graphics
% natbib is loaded once, above; the back-references ("Cited on page ...") in the
% bibliography come from hyperref's pagebackref option, which loads backref itself.
\usepackage{array}
\usepackage{latexsym}
\usepackage{etoolbox}
\usepackage{booktabs}
\usepackage{fancyvrb}

\usepackage{enumerate} % for numbering with [a)] format 

%----------------------------------------------------------------------------------------
% Maths
%----------------------------------------------------------------------------------------
\usepackage{xspace}

%----------------------------------------------------------------------------------------
% Own style for embedding python code in the text
%----------------------------------------------------------------------------------------
\usepackage{listings}
\usepackage{xcolor}

%----------------------------------------------------------------------------------------
% The xAILab palette (the lab's PowerPoint theme "xAILab Poster")
%----------------------------------------------------------------------------------------
% Defined here rather than in notation.tex because includes.tex is loaded first and
% both the listings style and hyperref below need them. The semantic commands that
% use them (\code, \highlightfile, \highlighttool, \highlightimportant, \blocklabel)
% live in notation.tex. Do not add colours outside this palette.
\definecolor{xaiInk}{HTML}{1C1C28}
\definecolor{xaiInkSoft}{HTML}{2B2B3A}
\definecolor{xaiTint}{HTML}{F3F3FD}
\definecolor{xaiAccent}{HTML}{5656E6}
\definecolor{xaiTool}{HTML}{814CB3}
\definecolor{xaiFile}{HTML}{B73F74}
\definecolor{xaiAlert}{HTML}{EE3334}
% literal colours for prose that names a colour in a figure (see \figgreen)
\definecolor{darkgreen}{rgb}{0,0.6,0}

\lstdefinestyle{mystyle}{
    backgroundcolor={},
    commentstyle=\color{xaiInkSoft},
    keywordstyle=\color{xaiAccent},
    stringstyle=\color{xaiFile},
    basicstyle=\ttfamily\footnotesize,
    breakatwhitespace=false,         
    breaklines=true,                 
    captionpos=b,                    
    keepspaces=true,               
    showspaces=false,                
    showstringspaces=false,
    showtabs=false,                  
    tabsize=2
}

\lstset{style=mystyle, escapeinside={(*@}{@*)}}

\newcommand\inlinecode[1]{{\lstinline!#1!}}



\ifxetex
\usepackage{fontspec}
\setmainfont[Scale=.8]{OpenDyslexic-Regular}
\else
\usepackage[pdftex,pagebackref,hypertexnames=false,colorlinks]{hyperref} % provide links in pdf
\hypersetup{pdftitle={},
  pdfsubject={}, 
  pdfauthor={},
  pdfkeywords={}, 
  pdfstartview=FitH,
  pdfpagemode={UseOutlines},% None, FullScreen, UseOutlines
  bookmarksnumbered=true, bookmarksopen=true, colorlinks=true,
    citecolor=xaiInk,%
    filecolor=xaiFile,%
    linkcolor=xaiAccent,%
    urlcolor=xaiAccent}
\usepackage[all]{hypcap}
\fi

\usepackage{tcolorbox}
\tcbuselibrary{skins,breakable}

%----------------------------------------------------------------------------------------
% The two boxed blocks
%----------------------------------------------------------------------------------------
% \begin{todo} ... \end{todo}        what the student has to implement
% \begin{attention} ... \end{attention}   a warning they must not miss
% Both are breakable, so a long list still splits across a page.
\newtcolorbox{todo}[1][]{%
  breakable, enhanced,
  colback=xaiTint, colframe=xaiAccent,
  boxrule=0pt, leftrule=2.5pt, arc=1pt, left=6pt, right=6pt, top=5pt, bottom=5pt,
  before upper={\textbf{\textcolor{xaiAccent}{ToDo}}\par\smallskip},
  #1}
\newtcolorbox{attention}[1][]{%
  breakable, enhanced,
  colback=white, colframe=xaiAlert,
  boxrule=0pt, leftrule=2.5pt, arc=1pt, left=6pt, right=6pt, top=5pt, bottom=5pt,
  before upper={\textbf{\textcolor{xaiAlert}{Attention}}\par\smallskip},
  #1}

%----------------------------------------------------------------------------------------
% One source, three documents
%----------------------------------------------------------------------------------------
% Every theory question used to exist three times: the question file, a _solution copy
% that reproduced the question verbatim and added the answer, and a third copy under
% student_version/ for the writeup template. They drifted, which is what item X1 was:
% a weight vector that was R^2 in the question and R^1 in the template, making the task
% unsolvable as stated. Now there is one file per question, rendered three ways.
%
%   \DLmode = 0   student description   answer hidden, space left for it
%           = 1   grader version        answer shown
%           = 2   writeup template      answer hidden, a fresh page for it
%
% main.tex does not set it. The build script selects a mode on the command line:
%
%   latexmk -pdf -usepretex='\def\DLmode{0}' main.tex
%
% and the default below is 1, so opening main.tex in Overleaf still gives the grader
% version, exactly as before.
\providecommand{\DLmode}{1}

\usepackage{comment}

% \begin{answer} ... \end{answer}  the worked solution. Only ever rendered in mode 1.
%
% IMPORTANT: this hides the answer, it does not remove it from the file. The writeup
% template shipped to students must therefore be built by a script that STRIPS these
% blocks (scripts/build_writeup_templates.sh), never by shipping the source with
% \DLmode=2. Hiding is a rendering decision; not shipping is a security one.
% Shown: an ordinary environment, so the body is typeset the way it always was.
% Hidden: \excludecomment from the comment package, which skips to \end{answer}
% without expanding anything, so an unfinished macro inside an answer cannot break
% the student description.
% (\specialcomment was tried for the shown case and is not worth it: it round-trips
% the body through a file and falls over on nested environments.)
% The label ends with \\ rather than \par on purpose. The old form was
% \teal{ Solution:\\ ... }, which leaves the paragraph open, and several answers rely
% on that: they put a \\ straight after a figure, which is legal inside a paragraph
% (a float does not break one) and an error in vertical mode. Ending the label with
% \par instead produced "There's no line here to end." in three answers.
\ifnum\DLmode=1
  \newenvironment{answer}%
    {\par\smallskip\noindent\color{xaiTool}\blocklabel{Solution:}\\}%
    {\par\smallskip}
\else
  \excludecomment{answer}
\fi

% \bymode{<description>}{<grader>}{<writeup>}
% The general form: emit different material in each of the three documents. Inside a
% question file it is now used only for grader-only spacing before an answer, e.g.
% \bymode{}{\vspace{0.5cm}}{}. Page breaks between questions use the two macros below.
\newcommand{\bymode}[3]{\ifcase\DLmode #1\or #2\or #3\fi}

% Page breaks between questions. These two macros are the ONLY separators a theory
% question file uses: no bare \newpage, \vfill or \bymode{...}{...}{\newpage} between
% questions. The layout they produce:
%   description and grader version  one numbered task (1.1, 1.2, ...) per page; its
%                                   lettered subtasks a), b), ... flow on that page
%   writeup template                one lettered subtask per page, because the student
%                                   writes the answer under it and the guidelines ask for
%                                   each problem/subproblem on a separate page
% Rules:
%   1. Always AFTER \end{answer}, never before \begin{answer}. The template strips the
%      answer and the student writes in its place; a break before it would put their
%      answer on the next question's page.
%   2. One macro per boundary: \subtaskbreak after a lettered subtask that another one
%      follows, \taskbreak after the last answer of a numbered task. A task without
%      lettered subtasks gets only \taskbreak. Lettered items that share one answer and
%      trailing notes (a4 4.2, 4.3) stay together.
%   3. theory.tex separates the question files with \taskbreak too.
% A \newpage inside a question's own text (a worked example before "Now it is your
% turn") or inside an answer is layout, not a separator, and is allowed.
\newcommand{\subtaskbreak}{\par\bymode{\medskip}{\medskip}{\newpage}}
\newcommand{\taskbreak}{\par\clearpage}


% various theorems
\usepackage{ntheorem}
\theoremstyle{break}
\newtheorem{lemma}{Lemma}
\newtheorem{theorem}{Theorem}
\newtheorem{remark}{Remark}
\newtheorem{definition}{Definition}
\newtheorem{proof}{Proof}

% example-environment
\newenvironment{example}[1][]
{ 
\vspace{4mm}
\noindent\makebox[\linewidth]{\rule{\hsize}{1.5pt}}
\textbf{Example #1}\\
}
{ 
\noindent\newline\makebox[\linewidth]{\rule{\hsize}{1.0pt}}
}

\ifxetex
\else
\renewcommand*{\rmdefault}{bch} % Charter
\renewcommand*{\ttdefault}{cmtt} % Computer Modern Typewriter
\fi

\setlength{\parindent}{0em}  % indentation of paragraph

\setlength{\headheight}{14.5pt}
\pagestyle{fancy}
\fancyfoot[ER,OL]{\thepage}%Page no. in the left on
                                %odd pages and on right on even pages
\fancyfoot[OC,EC]{\sffamily }

% One item in the running head, not two. fancyhdr's default puts \leftmark
% (section) and \rightmark (subsection) on the same line; these headings carry
% their point totals, so the two ran into each other on most pages of the coding
% parts. Show only the current section, and not in capitals.
\fancyhead{}
\fancyhead[LE,RO]{\slshape\nouppercase{\leftmark}}
\renewcommand{\sectionmark}[1]{\markboth{\thesection\quad #1}{}}
\renewcommand{\subsectionmark}[1]{}

\renewcommand{\headrulewidth}{0.1pt}
\renewcommand{\footrulewidth}{0.1pt}
\captionsetup{margin=10pt,font=small,labelfont=bf}


%--- chapter heading

\def\@makechapterhead#1{%
  \vspace*{10\p@}%
  {\parindent \z@ \raggedright
    \interlinepenalty\@M
    \Huge \bfseries 
    \thechapter \space\space #1\par\nobreak
    \vskip 30\p@
  }}

%---chapter heading for \chapter*  
\def\@makeschapterhead#1{%
  \vspace*{10\p@}%
  {\parindent \z@ \raggedright
    \sffamily
    \interlinepenalty\@M
    \Huge \bfseries  
    #1\par\nobreak
    \vskip 30\p@
  }}
  

% %%%%%%%%%%%%% boxit
\def\Beginboxit
   {\par
    \vbox\bgroup
	   \hrule
	   \hbox\bgroup
		  \vrule \kern1.2pt %
		  \vbox\bgroup\kern1.2pt
   }

\def\Endboxit{%
			      \kern1.2pt
		       \egroup
		  \kern1.2pt\vrule
		\egroup
	   \hrule
	 \egroup
   }	

\newenvironment{boxit}{\Beginboxit}{\Endboxit}
\newenvironment{boxit*}{\Beginboxit\hbox to\hsize{}}{\Endboxit}



\allowdisplaybreaks

\makeatletter
\newcounter{elimination@steps}
\newcolumntype{R}[1]{>{\raggedleft\arraybackslash$}p{#1}<{$}}
\def\elimination@num@rights{}
\def\elimination@num@variables{}
\def\elimination@col@width{}
\newenvironment{elimination}[4][0]
{
    \setcounter{elimination@steps}{0}
    \def\elimination@num@rights{#1}
    \def\elimination@num@variables{#2}
    \def\elimination@col@width{#3}
    \renewcommand{\arraystretch}{#4}
    \start@align\@ne\st@rredtrue\m@ne
}
{
    \endalign
    \ignorespacesafterend
}
\newcommand{\eliminationstep}[2]
{
    \ifnum\value{elimination@steps}>0\leadsto\quad\fi
    \left[
        \ifnum\elimination@num@rights>0
            \begin{array}
            {@{}*{\elimination@num@variables}{R{\elimination@col@width}}
            |@{}*{\elimination@num@rights}{R{\elimination@col@width}}}
        \else
            \begin{array}
            {@{}*{\elimination@num@variables}{R{\elimination@col@width}}}
        \fi
            #1
        \end{array}
    \right]
    & 
    \begin{array}{l}
        #2
    \end{array}
    &%                                    moved second & here
    \addtocounter{elimination@steps}{1}
}
\makeatother

%% Fast macro for column vectors
\makeatletter  
\def\colvec#1{\expandafter\colvec@i#1,,,,,,,,,\@nil}
\def\colvec@i#1,#2,#3,#4,#5,#6,#7,#8,#9\@nil{% 
  \ifx$#2$ \begin{bmatrix}#1\end{bmatrix} \else
    \ifx$#3$ \begin{bmatrix}#1\\#2\end{bmatrix} \else
      \ifx$#4$ \begin{bmatrix}#1\\#2\\#3\end{bmatrix}\else
        \ifx$#5$ \begin{bmatrix}#1\\#2\\#3\\#4\end{bmatrix}\else
          \ifx$#6$ \begin{bmatrix}#1\\#2\\#3\\#4\\#5\end{bmatrix}\else
            \ifx$#7$ \begin{bmatrix}#1\\#2\\#3\\#4\\#5\\#6\end{bmatrix}\else
              \ifx$#8$ \begin{bmatrix}#1\\#2\\#3\\#4\\#5\\#6\\#7\end{bmatrix}\else
                 \PackageError{Column Vector}{The vector you tried to write is too big, use bmatrix instead}{Try using the bmatrix environment}
              \fi
            \fi
          \fi
        \fi
      \fi
    \fi
  \fi 
}  
\makeatother

\robustify{\colvec}
